% === Begin Label Data ===
% ID: 450
% 难度星级: 
% 标签: 
% 备注: 
% 组卷引用次数: 0
% === End  Label Data ===

\begin{problem}{2017}{G}{北京卷}{14}{命题与逻辑}
三名工人加工同一种零件，他们在一天中的工作情况如图所示，其中点 $ A_i $ 的横、纵坐标分别为第 $ i $ 名工人上午的工作时间和加工的零件数，点 $ B_i $ 的横、纵坐标分别为第 $ i $ 名工人下午的工作时间和加工的零件数，$ i=1,2,3 $．
\begin{tikzpicture}[x=0.6cm,y=0.6cm,line cap=round,line join=round]
\coordinate (O) at (0,0);

% axes
\draw[->] (0,0) -- (8,0) node[below right] {工作时间(小时)};
\draw[->] (0,0) -- (0,5.5) node[above] {零件数(件)};
\node[below left] at (O) {$O$};

% points (approximate positions to match the figure)
\fill (3.2,4.4) circle (1.5pt) node[right] {$A_1$};
\fill (1.2,2.3) circle (1.5pt) node[above] {$A_2$};
\fill (2.6,0.7) circle (1.5pt) node[left] {$A_3$};

\fill (6.1,1.0) circle (1.5pt) node[right] {$B_1$};
\fill (5.4,2.8) circle (1.5pt) node[below] {$B_2$};
\fill (6.9,3.1) circle (1.5pt) node[right] {$B_3$};
\end{tikzpicture}
\circled{1} 记 $ Q_i $ 为第 $ i $ 名工人在这一天中加工的零件总数，则 $ Q_1,Q_2,Q_3 $ 中最大的是 \underline{\hspace{4em}}；

\circled{2} 记 $ p_i $ 为第 $ i $ 名工人在这一天中平均每小时加工的零件数，则 $ p_1,p_2,p_3 $ 中最大的是 \underline{\hspace{4em}}．
\end{problem}

\begin{answer}

\end{answer}

\begin{solutions}

\end{solutions}